\section{Limitations and Scope of Interpretation}
\label{sec:limitations}

\subsection{Purpose of this section}

This section brings together the limitations that affect the interpretation of
the study. They fall into three broad groups: limitations of the available
clinical data, limitations of the evaluation designs and limitations on how far
the results can be applied beyond the studied patients. These issues do not make
the analyses invalid, but they determine which conclusions the evidence can
support.

Table~\ref{tab:limitations:matrix} summarises the main limitations, their effect
on interpretation and the additional evidence needed to address them. The
proposed next steps describe future work; they were not part of the present
study.

\begin{table}[htbp]
\centering\scriptsize
\setlength{\tabcolsep}{3pt}
\caption{Main study limitations, their effect on interpretation and the
additional evidence needed to address them.}
\label{tab:limitations:matrix}
\begin{tabularx}{\linewidth}{@{}p{2.15cm}p{3.55cm}p{3.55cm}Y@{}}
\toprule
Area & Main limitation & Effect on interpretation & What future work should add \\
\midrule
\csname @@input\endcsname{tables/limitations/limitations_matrix_rows.tex}
\bottomrule
\end{tabularx}
\end{table}

\subsection{Data source, cohort selection and measurement}

The study is retrospective and uses one release of the Surveillance,
Epidemiology, and End Results (SEER) programme. The diagnostic cohort contains
2,672 patients after applying the eligibility and preprocessing procedure based
on the reference study. This represents 46.4\% of the 5,761 exported records.
The largest exclusion removed 1,903 records with at least one unknown required
field. Because the analysis uses complete records, the retained patients may
differ systematically from those who were excluded. The models describe the
available cohort; they cannot correct for patients or clinical information that
are absent from it.

Several variables are simplified versions of the underlying clinical
information. Stage values from different diagnosis years and assessment types
were combined into common categories. Histology and pancreatic site were also
grouped, age was top-coded, and some treatment codes combine ``no treatment''
with ``unknown''. The data do not establish whether every treatment variable
would have been known at the intended prediction time. Associations involving
treatment should therefore be interpreted as predictive relationships, not as
evidence that a treatment caused an outcome. A model intended for prospective
use would need clearly time-stamped predictors available before the prediction
is made.

The outcomes are also specific. The diagnostic task predicts liver metastasis
recorded at the initial diagnosis; it does not predict metastasis that develops
later. The prognostic task evaluates all-cause survival in the 462 patients who
already had liver metastasis at diagnosis. It does not estimate cancer-specific
survival and does not represent all patients with pancreatic neuroendocrine
tumours. Twelve prognostic records have zero recorded survival months. They were
retained as recorded, with a positive-time adjustment used only where required
by a model. This convention permits the analyses to run but cannot resolve
possible uncertainty in the recorded timing.

The SEER release used here is newer than the dataset used in the reference
study, because the earlier release was no longer available. Differences between
the published and reproduced results may therefore arise from changes in the
patients, calendar period, variable definitions or implementation. The work is
best understood as a reconstruction of the published methodology on the
available SEER data, rather than an exact replication on the original cohort.

\subsection{External validation and statistical uncertainty}

The reproduced reference models were evaluated without modification in the
independent University Hospital Erlangen (UK Erlangen) cohort. This is an
important external test, but the cohort is small. It contains 30 patients for
the diagnostic task and 11 patients with liver metastasis for the prognostic
task. The four diagnostic models could be compared in a common set of 29
patients because one patient had missing nodal stage. The prognostic nomogram
could be evaluated in all 11 patients, whereas its comparison with the tumour-
and nodal-stage model used 10. Confidence intervals were consequently wide,
survival-risk groups contained only three to seven patients, and the prognostic
sample was too small for decision-curve analysis. These results show whether the
unchanged models can be applied to the available hospital data, but they cannot
establish stable clinical thresholds, subgroup safety or performance across
different hospitals.

The foundation-model, multilayer perceptron (MLP), embedding and synthetic-data
experiments were evaluated only within SEER. Repeated partitions reduce
dependence on one arbitrary train--test division, but all partitions still come
from the same registry release. They do not replace evaluation in patients from
a different time period, institution or country.

The report also contains several evaluation designs that answer different
questions. The reference-study reconstruction, external evaluation, complete-
cohort embedding analyses, synthetic-cohort analyses and repeated held-out
SEER comparisons do not use the same patients or the same validation procedure.
Their performance values should therefore be interpreted within their own
designs and should not be combined into a single ranking of all models.

Most 95\% confidence intervals were calculated by resampling the evaluated
patients while keeping the fitted models unchanged. They quantify uncertainty
associated with the available evaluation patients, but not the additional
variation that could arise from repeating feature selection, model fitting,
representation learning, synthetic-data generation or foundation-model
pretraining. Most intervals are also pointwise and were not adjusted for the
number of models, metrics and augmentation sizes examined. An interval that
excludes zero supports a difference for that particular comparison; it does not
by itself establish clinical superiority. Likewise, an interval that includes
zero does not prove that two methods are equivalent. The three MLP
initialisations provide a limited sensitivity check rather than a complete
description of optimisation variability.

Some secondary results have less statistical support. The planned confidence
intervals for the prognostic Kaplan--Meier calibration comparisons in the
augmentation experiment could not be reported because too few bootstrap
samples retained adequate numbers of patients at risk. Decision-curve results
were descriptive and no risk threshold was prospectively selected for clinical
action. The study therefore does not demonstrate that using any model would
improve patient outcomes.

\subsection{Learned representations and synthetic data}

The Tabular Prior-Data Fitted Network (TabPFN) embeddings were created from each
complete task cohort using outcome information. A patient's own outcome was
excluded when its representation was generated, but the transformation itself
was not rebuilt separately within each later training partition. The embedding
analyses show that the representations retain predictive information, but they
do not establish how an embedding learned only from training patients would
perform in previously unseen patients. The prognostic time representation also
used observed follow-up time as a regression target and did not model censoring
directly, which further limits its interpretation as a survival representation.

The synthetic cohorts generated from the complete SEER datasets have the same
dependency on the full observed cohort. They are useful for studying whether
distributions and modelling pipelines can be reproduced, but they are not
independent validation data. In the augmentation experiment, this dependency
was reduced by training a new generator only on the real training patients in
each evaluation partition and testing the models only on held-out real
patients. Even under this stronger design, generated records do not add new
clinical observations. At the largest prognostic augmentation size, synthetic
records formed 97.9\% of the combined training data, so the fitted model was
dominated by the generator's approximation of a small real cohort.

The fidelity results identify clinically relevant differences between real and
synthetic records. Grade 4 was absent from every generated dataset, some
continuous tumour sizes disagreed with their recorded size groups, and outcome
frequency, survival time and relationships between variables were not always
preserved. Similar distributions for individual variables therefore do not
guarantee preservation of the combinations that matter to a prediction model.
Conversely, the improvement observed for diagnostic XGBoost does not establish
that the synthetic data were faithful in every respect.

No formal privacy analysis was performed. The absence of exact duplicate rows
and the use of new identifiers are useful integrity checks, but they do not
measure whether information about individual patients could be inferred from
the generated data. The synthetic records should therefore not be considered
anonymous clinical substitutes or released on the basis of this study alone.

\subsection{Model scope, subgroups and clinical applicability}

The evaluated model set is substantial but not exhaustive. The MLP experiment
tested one prespecified architecture and training procedure; it was not a search
over all possible neural networks. The foundation-model results apply to the
specific model versions, preprocessing choices and inference procedures used in
this study. Their pretraining datasets were not independently reconstructed, so
the analysis cannot determine how much of the observed performance arose from
pretraining or rule out overlap with similar data sources. The absence of a
prognostic advantage in this study does not prove that all future
survival-focused foundation models will perform similarly.

Rare clinical groups remain particularly uncertain. The prognostic cohort
contains only three patients with brain metastasis and 33 with grade 4. The
diagnostic cohort contains 60 patients with grade 4, but many individual
evaluation partitions had too few such patients for a stable estimate. No
prognostic partition met the prespecified support requirement for grade-4
performance, and synthetic augmentation could not improve this because grade 4
was absent from the generated data. A missing subgroup estimate is not evidence
of equal performance, and strong overall results do not establish fairness or
safety in rare groups.

Finally, the study is predictive rather than causal. It does not estimate the
effect of treatment, define a clinical intervention, compare healthcare
workflows or measure prospective patient benefit. Discrimination, calibration,
prediction error and decision curves are useful components of model evaluation,
but they are not sufficient for deployment. Clinical use would require larger
external cohorts, validation across time and centres, predefined decision
thresholds, assessment of rare groups and a prospective study of clinical
impact.

\subsection{Implications for the conclusions}

Within these limits, the matched internal comparisons support small diagnostic
advantages for the three evaluated tabular foundation models over logistic
regression. They also support a specific diagnostic augmentation benefit for
XGBoost. The analyses did not establish a general advantage for the tested MLP,
synthetic augmentation or foundation models in prognosis. The independent
Erlangen results provide encouraging evidence that the reproduced reference
models retained predictive ranking, but the sample is too small to determine
their performance precisely.

The study does not establish external performance for the foundation-model and
synthetic-data experiments, causal benefit, formal privacy, reliable performance
in rare subgroups or readiness for clinical deployment. Future work should
evaluate all model families in larger and temporally distinct multicentre
cohorts, learn every representation and synthetic-data generator using training
patients only, recruit enough patients from rare clinical groups, repeat the
complete model-development process when estimating uncertainty and assess both
privacy risk and prospective clinical benefit.
